% Part: incompleteness
% Chapter: representability-in-q
% Section: basic-representable

\documentclass[../../../include/open-logic-section]{subfiles}

\begin{document}

\olfileid{inc}{req}{bre}
\olsection{মৌলিক অপেক্ষকগুলি $\Th{Q}$-তে উপস্থাপনযোগ্য}

প্রথমে দেখাতে হবে যে সব মৌলিক অপেক্ষক $\Th{Q}$-তে উপস্থাপনযোগ্য। শেষ পর্যন্ত
প্রত্যেক $k$-স্থানীয় মৌলিক অপেক্ষক $f(x_0,\dots,x_{k-1})$-কে এমন একটি
!!a{formula} $!A_f(x_0,\dots,x_{k-1},y)$ দেওয়ার উপায় দেখাতে হবে, যা
অপেক্ষকটিকে উপস্থাপন করে।

শূন্য, উত্তরসূরি, যোগ, গুণ, সমতার চরিত্রসূচক অপেক্ষক এবং অভিক্ষেপগুলিকে
উপস্থাপন করা যাবে। প্রতিটি ক্ষেত্রেই উপযুক্ত উপস্থাপক সূত্র একেবারেই সরল;
যেমন, শূন্যকে $y = \Obj 0$ সূত্রটি, উত্তরসূরিকে !!{formula} $x_0' = y$ এবং
যোগকে !!{formula} $(x_0 + x_1) = y$ উপস্থাপন করে। আসল কাজ হলো দেখানো যে
$\Th{Q}$ প্রাসঙ্গিক !!{sentence}s প্রমাণ করতে পারে; যেমন, উপরের !!{formula}
যোগকে উপস্থাপন করে বলার জন্য দেখাতে হয় যে প্রত্যেক জোড়া স্বাভাবিক সংখ্যা
$m$ ও $n$-এর জন্য $\Th{Q}$ প্রমাণ করে
\begin{align*}
& \eq[\num n + \num m][\num {n+m}] \text{ এবং}\\
& \lforall[y][(\eq[(\num n + \num m)][y] \lif \eq[y][\num{n+m}])].
\end{align*}

\begin{prop}
\ollabel{prop:rep-zero}
শূন্য অপেক্ষক $\Zero(x) = 0$-কে $\Th{Q}$-তে
$!A_{\Zero}(x,y) \ident \eq[y][\Obj 0]$ উপস্থাপন করে।
\end{prop}

\begin{prop}
\ollabel{prop:rep-succ}
উত্তরসূরি অপেক্ষক $\Succ(x) = x+1$-কে $\Th{Q}$-তে
$!A_{\Succ}(x,y) \ident \eq[y][x']$ উপস্থাপন করে।
\end{prop}

\begin{prop}
\ollabel{prop:rep-proj}
অভিক্ষেপ অপেক্ষক $\Proj{n}{i}(x_0, \dots, x_{n-1}) = x_i$-কে
$\Th{Q}$-তে
\[!A_{\Proj{n}{i}}(x_0, \dots, x_{n-1}, y) \ident \eq[y][x_i].\]
উপস্থাপন করে।
\end{prop}

\begin{prob}
প্রমাণ করো যে $\eq[y][\Obj 0]$, $\eq[y][x']$ ও $\eq[y][x_i]$ যথাক্রমে
$\Zero$, $\Succ$ ও $\Proj{n}{i}$-কে উপস্থাপন করে।
\end{prob}

\begin{prop}
\ollabel{prop:rep-id}
$=$-এর চরিত্রসূচক অপেক্ষক
\[
\Char{=}(x_0, x_1) =
\begin{cases}
  1 & \text{যদি } x_0 =x_1\\
  0 & \text{অন্যথায়}
\end{cases}
\]
% BN-SRC-233 TARGET-CORRECTION-BEGIN
\par\smallskip\noindent\textnormal{\small\textbf{উৎস-সংশোধন (BN-SRC-233):} হিমায়িত \texttt{cases} বিন্যাসে প্রথম শর্তটি \texttt{\string\text\{...\}}-এর মধ্যে থাকলেও দ্বিতীয় পংক্তির ``otherwise'' সরাসরি গণিত-মোডে ছিল, ফলে সেটি শব্দের বদলে চলরাশির গুণফলের মতো মুদ্রিত হতো। বাংলা পাঠে দ্বিতীয় পংক্তির পাঠও \texttt{\string\text\{...\}}-এর মধ্যে রাখা হয়েছে।} % BN-SRC-233 TARGET-CORRECTION-END
$\Th{Q}$-তে নিচের সূত্র দ্বারা উপস্থাপিত:
\[
  !A_{\Char{=}}(x_0, x_1, y) \ident (\eq[x_0][x_1] \land \eq[y][\num{1}]) \lor (\eq/[x_0][x_1] \land
\eq[y][\num{0}]).
\]
\end{prop}

প্রমাণটির জন্য নিচের সহায়ক উপপাদ্য দরকার।

\begin{lem}
\ollabel{lem:q-proves-neq} স্বাভাবিক সংখ্যা $n$ ও $m$ দেওয়া থাকলে, $n
\neq m$ হলে $\Th{Q} \Proves \eq/[\num n][\num m]$।
\end{lem}

\begin{proof}
$n$-এর উপর আরোহ ব্যবহার করে দেখাও: প্রত্যেক $m$-এর জন্য $n \neq m$ হলে
$Q \Proves \eq/[\num n][\num m]$।

ভিত্তি ক্ষেত্রে $n = 0$। $m$, $0$-র সমান না হলে কোনো স্বাভাবিক সংখ্যা
$k$-এর জন্য $m = k + 1$। আমাদের একটি স্বতঃসিদ্ধ বলে
$\lforall[x][\eq/[0][x']]$। একটি পরিমাণসূচক-স্বতঃসিদ্ধে $x$-এর স্থানে
$\num k$ বসিয়ে $\eq/[0][\num k']$ সিদ্ধান্তে আসতে পারি। কিন্তু
$\num k'$-ই $\num m$।

আরোহের ধাপে দাবিটি $n$-এর জন্য সত্য ধরে $n+1$ বিবেচনা করি। $m$ যেকোনো
স্বাভাবিক সংখ্যা হোক। দুটি সম্ভাবনা: হয় $m = 0$, নয়তো কোনো $k$-এর জন্য
$m = k+1$। প্রথম ক্ষেত্রটি উপরের মতোই সামলানো যায়। দ্বিতীয় ক্ষেত্রে ধরি
$n+1 \neq k+1$। তাহলে $n \neq k$। $n$-এর জন্য আরোহের অনুমান থেকে পাই
$\Th{Q} \Proves \eq/[\num n][\num k]$। একটি স্বতঃসিদ্ধ বলে
$\lforall[x][\lforall[y][\eq[x'][y'] \lif \eq[x][y]]]$। একটি
পরিমাণসূচক-স্বতঃসিদ্ধ ব্যবহার করে পাই
$\eq[\num n'][\num k'] \lif \eq[\num n][\num k]$। বচনগত যুক্তিবিদ্যা
ব্যবহার করে $\Th{Q}$-তে সিদ্ধান্তে আসতে পারি
$\eq/[\num n][\num k] \lif \eq/[\num n'][\num k']$। modus ponens
ব্যবহার করে পাই $\eq/[\num n'][\num k']$; এটিই আমাদের দরকার, কারণ
$\num k'$ হলো $\num m$।
\end{proof}

\begin{explain}
সহায়ক উপপাদ্যটি যে খুব বেশি কিছু বলে তা নয়: মূলত বলে, আলাদা সংখ্যাপদ
আলাদা বস্তুকে নির্দেশ করে—এ কথা $\Th{Q}$ প্রমাণ করতে পারে। যেমন, $\Th{Q}$
$0'' \neq 0'''$ প্রমাণ করে। কিন্তু সাধারণভাবে কথাটি দেখাতে কিছু সতর্কতা
দরকার। আরও লক্ষ করো, আমরা আরোহ ব্যবহার করলেও সেটি $\Th{Q}$-এর
\emph{বাইরের} আরোহ।
\end{explain}

\begin{proof}[\olref{prop:rep-id}-এর প্রমাণ]
$n = m$ হলে $\num{n}$ ও $\num{m}$ একই পদ এবং $\Char{=}(n, m) = 1$। কিন্তু
$\Th{Q} \Proves (\eq[\num{n}][\num{m}] \land
\eq[\num{1}][\num{1}])$, তাই এটি $!A_=(\num{n}, \num{m},
\num{1})$ প্রমাণ করে। $n \neq m$ হলে $\Char=(n, m) = 0$।
\olref{lem:q-proves-neq} অনুসারে $\Th{Q} \Proves
\eq/[\num{n}][\num{m}]$, এবং তাই
$(\eq/[\num{n}][\num{m}] \land \Obj 0 = \Obj 0)$-ও। অতএব
$\Th{Q} \Proves !A_=(\num{n}, \num{m}, \num{0})$।

দ্বিতীয় অংশের জন্যও দুটি ক্ষেত্র আছে। $n = m$ হলে দেখাতে হবে
$\Th{Q} \Proves \lforall[y][(!A_=(\num{n}, \num{m}, y)
  \lif \eq[y][\num{1}])]$। অনানুষ্ঠানিকভাবে যুক্তি করি: ধরি
$!A_=(\num{n}, \num{m}, y)$, অর্থাৎ
\[
(\eq[\num{n}][\num{n}] \land \eq[y][\num{1}]) \lor
(\eq/[\num{n}][\num{n}] \land \eq[y][\num{0}])
\]
বাম বিযোজ্যটি যুক্তিবিদ্যা থেকে $\eq[y][\num{1}]$ দেয়; ডান বিযোজ্যটি
যুক্তিবিদ্যায় প্রমাণযোগ্য $\eq[\num{n}][\num{n}]$-এর বিরোধী।

অন্য দিকে, ধরি $n \neq m$। তখন $!A_=(\num{n}, \num{m}, y)$ হলো
\[
(\eq[\num{n}][\num{m}] \land \eq[y][\num{1}]) \lor
(\eq/[\num{n}][\num{m}] \land \eq[y][\num{0}])
\]
এখানে বাম বিযোজ্যটি $\eq/[\num{n}][\num{m}]$-এর বিরোধী; শেষের সূত্রটি
\olref{lem:q-proves-neq} অনুসারে $\Th{Q}$-তে প্রমাণযোগ্য। ডান বিযোজ্য থেকে
$\eq[y][\num{0}]$ পাওয়া যায়।
\end{proof}

\begin{prop}
\ollabel{prop:rep-add}
যোগ অপেক্ষক $\Add(x_0, x_1) = x_0+x_1$-কে $\Th{Q}$-তে উপস্থাপন করে
\[
  !A_{\Add}(x_0, x_1, y) \ident \eq[y][(x_0 + x_1)].
\]
\end{prop}

\begin{lem}
\ollabel{lem:q-proves-add}
$\Th{Q} \Proves \eq[(\num{n} + \num{m})][\num{n+m}]$
\end{lem}

\begin{proof}
$m$-এর উপর আরোহ দিয়ে এটি প্রমাণ করি। $m = 0$ হলে দাবিটি হলো
$\Th{Q} \Proves \eq[(\num{n} + \Obj 0)][\num{n}]$। এটি স্বতঃসিদ্ধ~$!Q_4$
থেকে আসে। এবার দাবিটি $m$-এর জন্য ধরে $m+1$-এর জন্য প্রমাণ করি; অর্থাৎ
$\Th{Q} \Proves \eq[(\num{n} + \num{m+1})][\num{n+m+1}]$ প্রমাণ করি।
লক্ষ করো, $\num{m+1}$ হলো $\num{m}'$, আর $\num{n+m+1}$ হলো
$\num{n+m}'$। স্বতঃসিদ্ধ $!Q_5$ অনুসারে
$\Th{Q} \Proves \eq[(\num{n} + \num{m}')][(\num{n}+\num{m})']$। আরোহের
অনুমান অনুসারে $\Th{Q} \Proves
\eq[(\num{n} + \num{m})][\num{n+m}]$। তাই
$\Th{Q} \Proves \eq[(\num{n} + \num{m}')][\num{n+m}']$।
\end{proof}

\begin{proof}[\olref{prop:rep-add}-এর প্রমাণ]
$\Add$-কে উপস্থাপন করা !!{formula}~$!A_\Add(x_0, x_1, y)$ হলো
$\eq[y][(x_0 + x_1)]$। প্রথমে দেখাই, $\Add(n, m) = k$ হলে
$\Th{Q} \Proves !A_\Add(\num{n}, \num{m}, \num{k})$; অর্থাৎ
$\Th{Q} \Proves \eq[\num{k}][(\num{n} + \num{m})]$। কিন্তু
$k = n + m$ বলে $\num{k}$-ই $\num{n+m}$, এবং
\olref{lem:q-proves-add}-এ দেখানো হয়েছে যে
$\Th{Q} \Proves \eq[(\num{n} + \num{m})][\num{n+m}]$।

আরও দেখাতে হবে, $\Add(n, m) = k$ হলে
\[
\Th{Q} \Proves \lforall[y][(!A_\Add(\num{n}, \num{m}, y) \lif
  \eq[y][\num{k}])].
\]
ধরি $\eq[(\num{n} + \num{m})][y]$। যেহেতু
\[
\Th{Q} \Proves \eq[(\num{n}+\num{m})][\num{n+m}],
\]
তাই বাম দিকটির স্থানে $\num{n+m}$ বসিয়ে যেকোনো~$y$-এর জন্য
$\eq[\num{n+m}][y]$ পাই।
\end{proof}

\begin{prop}
\ollabel{prop:rep-mult}
গুণ অপেক্ষক $\Mult(x_0, x_1) = x_0 \cdot x_1$-কে $\Th{Q}$-তে উপস্থাপন করে
\[
  !A_{\Mult}(x_0, x_1, y) \ident y = (x_0 \times x_1).
\]
\end{prop}

\begin{proof} অনুশীলন। \end{proof}

\begin{lem}
\ollabel{lem:q-proves-mult}
$\Th{Q} \Proves \eq[(\num{n} \times \num{m})][\num{n \cdot m}]$
\end{lem}

\begin{proof} অনুশীলন। \end{proof}

\begin{prob}
\olref[inc][req][bre]{lem:q-proves-mult} প্রমাণ করো।
\end{prob}

\begin{prob}
\olref[inc][req][bre]{lem:q-proves-mult} ব্যবহার করে
\olref[inc][req][bre]{prop:rep-mult} প্রমাণ করো।
\end{prob}

\begin{explain}
  মনে করো, পাটিগণিতের ভাষার অপেক্ষক প্রতীকের জন্য আমরা $\times$ এবং সংখ্যার
  উপর সাধারণ গুণ ক্রিয়ার জন্য $\cdot$ ব্যবহার করি। তাই $\cdot$ সংখ্যার
  রাশির মধ্যে বসতে পারে (যেমন $m \cdot n$-এ), আর $\times$ শুধু পাটিগণিতের
  ভাষার পদের মধ্যে বসে (যেমন $(\num{m} \times \num{n})$-এ)। আরও বিভ্রান্তিকর
  হলো, !!{function} ও যোগ ক্রিয়া—দুটির জন্যই $+$ ব্যবহার করা হয়। পদের মধ্যে
  থাকলে—যেমন $(\num{n} + \num{m})$-এ—এটি পাটিগণিতের ভাষার $2$-স্থানীয়
  !!{function}; আর সংখ্যার মধ্যে থাকলে—যেমন $n+m$-এ—এটি যোগ ক্রিয়া। এর
  মধ্যে $\num{n+m}$-ও পড়ে: এটি সংখ্যা~$n+m$-এর সংশ্লিষ্ট প্রমিত সংখ্যাপদ।
\end{explain}

\end{document}
